% GENERATED FILE. Edit canonical lessons, then run npm run generate:books.
% canonical-source-hash: fnv1a64:c7761520a4442103
% canonical-lessons: UR-C331-sabab, UR-C331-bunyad, UR-C331-zimme-dar, UR-C331-jurna, UR-C331-vabasta

\chapter{What Caused It}
\label{ch:ur-what-caused-it}
\hlchaptermodality{331}

\begin{chapteropening}
\textbf{By the end of this chapter:} \emph{I can say what caused something and who was responsible.}

Complete the last lesson of chapter 331: I can say what caused something and who was responsible.
\end{chapteropening}

\section[\texorpdfstring{\ur{سبب} (sabab)}{سبب (sabab)}]{\ur{سبب} (sabab) — a cause}

\label{lesson:UR-C331-sabab}



\begin{quote}
Before the new one: say the Urdu for a fruit seller, then the Urdu for a greengrocer.
\end{quote}

\subsection*{You'll want to know: \ur{سبب}}
\textbf{\ur{سبب}} — \emph{sabab} — \textquotedblleft{}a cause\textquotedblright{}.

\begin{quote}
\ur{اس} \ur{حادثے} \ur{کا} \ur{سبب} \ur{کیا} \ur{تھا؟}
\end{quote}

\begin{quote}
\emph{is hādse kā sabab kyā thā?}
\end{quote}

\begin{quote}
What was the cause of this accident?
\end{quote}

Masculine.

\begin{cousinweb}
\textbf{\ur{سبب}} (\emph{sabab}) is usually traced, through Persian, to Arabic \emph{sabab}, \textquotedblleft{}a cause, a means\textquotedblright{}.
\end{cousinweb}

\begin{grammarlens}[title={What you've built}]
One more word for giving reasons and opinions.
\end{grammarlens}

\subsection*{Your turn}
\begin{itemize}
\raggedright
  \item \emph{Say it:} \emph{sabab}
  \item \emph{Say it:} \emph{sabab}, once more
  \item \emph{Say it:} \emph{is hādse kā sabab kyā thā?} — What was the cause of this accident?
  \item \emph{Recall:} say the Urdu for a fruit seller, then the Urdu for a greengrocer, then say \emph{sabab} again
\end{itemize}

\begin{tcolorbox}[breakable,colback=teal!4,colframe=teal!35!black,title={Before you move on}]
What does \ur{سبب} mean? (\textquotedblleft{}A cause\textquotedblright{}.) Say it once more.
\end{tcolorbox}
\section[\texorpdfstring{\ur{بنیاد} (bunyād)}{بنیاد (bunyād)}]{\ur{بنیاد} (bunyād) — a basis; a foundation}

\label{lesson:UR-C331-bunyad}



\begin{quote}
Before the new one: say the Urdu for a greengrocer, then the Urdu for a cause.
\end{quote}

\subsection*{You'll want to know: \ur{بنیاد}}
\textbf{\ur{بنیاد}} — \emph{bunyād} — \textquotedblleft{}a basis; a foundation\textquotedblright{}.

\begin{quote}
\ur{آپ} \ur{کی} \ur{بات} \ur{کی} \ur{بنیاد} \ur{کیا} \ur{ہے؟}
\end{quote}

\begin{quote}
\emph{āp kī bāt kī bunyād kyā hai?}
\end{quote}

\begin{quote}
What is your point based on?
\end{quote}

Feminine.

\begin{cousinweb}
\textbf{\ur{بنیاد}} (\emph{bunyād}) is usually traced to Persian \emph{bunyād}, \textquotedblleft{}a foundation\textquotedblright{}; Urdu uses it for the ground an argument stands on as well as for the base of a building.
\end{cousinweb}

\begin{grammarlens}[title={What you've built}]
One more word for giving reasons and opinions.
\end{grammarlens}

\subsection*{Your turn}
\begin{itemize}
\raggedright
  \item \emph{Say it:} \emph{bunyād}
  \item \emph{Say it:} \emph{bunyād}, once more
  \item \emph{Say it:} \emph{āp kī bāt kī bunyād kyā hai?} — What is your point based on?
  \item \emph{Recall:} say the Urdu for a greengrocer, then the Urdu for a cause, then say \emph{bunyād} again
\end{itemize}

\begin{tcolorbox}[breakable,colback=teal!4,colframe=teal!35!black,title={Before you move on}]
What does \ur{بنیاد} mean? (\textquotedblleft{}A basis; a foundation\textquotedblright{}.) Say it once more.
\end{tcolorbox}
\section[\texorpdfstring{\ur{ذمہ} \ur{دار} (zimme-dār)}{ذمہ دار (zimme-dār)}]{\ur{ذمہ} \ur{دار} (zimme-dār) — responsible}

\label{lesson:UR-C331-zimme-dar}



\begin{quote}
Before the new one: say the Urdu for a cause, then the Urdu for a basis.
\end{quote}

\subsection*{You'll want to know: \ur{ذمہ} \ur{دار}}
\textbf{\ur{ذمہ} \ur{دار}} — \emph{zimme-dār} — \textquotedblleft{}responsible\textquotedblright{}.

\begin{quote}
\ur{سارا} \ur{اپنے} \ur{کام} \ur{کی} \ur{ذمہ} \ur{دار} \ur{ہے۔}
\end{quote}

\begin{quote}
\emph{sārā apne kām kī zimme-dār hai.}
\end{quote}

\begin{quote}
Sara is responsible for her own work.
\end{quote}

\begin{cousinweb}
\textbf{\ur{ذمہ} \ur{دار}} (\emph{zimme-dār}) is usually traced to Persian \emph{zimmadār}, \textquotedblleft{}answerable, in charge\textquotedblright{}: \emph{zimma}, a charge one carries, with \emph{-dār}, \textquotedblleft{}holding\textquotedblright{}.
\end{cousinweb}

\begin{grammarlens}[title={What you've built}]
One more word for giving reasons and opinions.
\end{grammarlens}

\subsection*{Your turn}
\begin{itemize}
\raggedright
  \item \emph{Say it:} \emph{zimme-dār}
  \item \emph{Say it:} \emph{zimme-dār}, once more
  \item \emph{Say it:} \emph{sārā apne kām kī zimme-dār hai.} — Sara is responsible for her own work.
  \item \emph{Recall:} say the Urdu for a cause, then the Urdu for a basis, then say \emph{zimme-dār} again
\end{itemize}

\begin{tcolorbox}[breakable,colback=teal!4,colframe=teal!35!black,title={Before you move on}]
What does \ur{ذمہ} \ur{دار} mean? (\textquotedblleft{}Responsible\textquotedblright{}.) Say it once more.
\end{tcolorbox}
\section[\texorpdfstring{\ur{جڑنا} (juṛnā)}{جڑنا (juṛnā)}]{\ur{جڑنا} (juṛnā) — to be joined; to be linked}

\label{lesson:UR-C331-jurna}



\begin{quote}
Before the new one: say the Urdu for a basis, then the Urdu for responsible.
\end{quote}

\subsection*{You'll want to know: \ur{جڑنا}}
\textbf{\ur{جڑنا}} — \emph{juṛnā} — \textquotedblleft{}to be joined; to be linked\textquotedblright{}.

\begin{quote}
\ur{یہ} \ur{سڑک} \ur{اس} \ur{پل} \ur{سے} \ur{جڑتی} \ur{ہے۔}
\end{quote}

\begin{quote}
\emph{yih saṛak is pul se juṛtī hai.}
\end{quote}

\begin{quote}
This road links up with this bridge.
\end{quote}

Past (it …): \ur{وہ} \ur{پل} \ur{سڑک} \ur{سے} \ur{جڑ} \ur{گیا}, \emph{voh pul saṛak se juṛ gayā}.

\begin{cousinweb}
\textbf{\ur{جڑنا}} (\emph{juṛnā}) is inherited: CDIAL lists the Hindi cousin \emph{juṛnā}, \textquotedblleft{}to be joined\textquotedblright{}, under \emph{*yuṭati}, a form it reconstructs from its descendants (CDIAL 10496).
\end{cousinweb}

\begin{grammarlens}[title={What you've built}]
One more word for giving reasons and opinions.
\end{grammarlens}

\subsection*{Your turn}
\begin{itemize}
\raggedright
  \item \emph{Say it:} \emph{juṛnā}
  \item \emph{Say it:} \emph{juṛnā}, once more
  \item \emph{Say it:} \emph{yih saṛak is pul se juṛtī hai.} — This road links up with this bridge.
  \item \emph{Recall:} say the Urdu for a basis, then the Urdu for responsible, then say \emph{juṛnā} again
\end{itemize}

\begin{tcolorbox}[breakable,colback=teal!4,colframe=teal!35!black,title={Before you move on}]
What does \ur{جڑنا} mean? (\textquotedblleft{}To be joined; to be linked\textquotedblright{}.) Say it once more.
\end{tcolorbox}
\section[\texorpdfstring{\ur{وابستہ} (vābastā)}{وابستہ (vābastā)}]{\ur{وابستہ} (vābastā) — connected (with); tied (to)}

\label{lesson:UR-C331-vabasta}



\begin{quote}
Before the new one: say the Urdu for responsible, then the Urdu for to be linked.
\end{quote}

\subsection*{You'll want to know: \ur{وابستہ}}
\textbf{\ur{وابستہ}} — \emph{vābastā} — \textquotedblleft{}connected (with); tied (to)\textquotedblright{}.

\begin{quote}
\ur{اس} \ur{کا} \ur{نام} \ur{اس} \ur{شہر} \ur{سے} \ur{وابستہ} \ur{ہے۔}
\end{quote}

\begin{quote}
\emph{us kā nām is shahr se vābastā hai.}
\end{quote}

\begin{quote}
His name is bound up with this city.
\end{quote}

\begin{cousinweb}
\textbf{\ur{وابستہ}} (\emph{vābastā}) is usually traced to Persian \emph{vābasta}, \textquotedblleft{}tied back, attached\textquotedblright{}, built on \emph{basta}, \textquotedblleft{}tied\textquotedblright{}.
\end{cousinweb}

\begin{grammarlens}[title={What you've built}]
One more word for giving reasons and opinions.
\end{grammarlens}

\subsection*{Your turn}
\begin{itemize}
\raggedright
  \item \emph{Say it:} \emph{vābastā}
  \item \emph{Say it:} \emph{vābastā}, once more
  \item \emph{Say it:} \emph{us kā nām is shahr se vābastā hai.} — His name is bound up with this city.
  \item \emph{Recall:} say the Urdu for responsible, then the Urdu for to be linked, then say \emph{vābastā} again
\end{itemize}

\begin{tcolorbox}[breakable,colback=teal!4,colframe=teal!35!black,title={Before you move on}]
What does \ur{وابستہ} mean? (\textquotedblleft{}Connected (with); tied (to)\textquotedblright{}.) Say it once more.
\end{tcolorbox}
